\documentclass[conference]{IEEEtran}
\IEEEoverridecommandlockouts
\usepackage{cite}
\usepackage{amsmath,amssymb,amsfonts}
\usepackage{algorithmic}
\usepackage{graphicx}
\usepackage{textcomp}
\usepackage{xcolor}
\usepackage{hyperref}
\newcommand{\srclink}[2]{\href{#1}{\textcolor{blue}{#2}}}
\usepackage[ruled,vlined]{algorithm2e}
\usepackage{booktabs}
\usepackage{multirow}

\def\BibTeX{{\rm B\kern-.05em{\sc i\kern-.025em b}\kern-.08em
    T\kern-.1667em\lower.7ex\hbox{E}\kern-.125emX}}

\setlength{\parindent}{0pt}
\usepackage{caption}
\setlength{\textfloatsep}{6pt plus 1.0pt minus 1.5pt}
\setlength{\floatsep}{6pt plus 1.0pt minus 1.5pt}
\setlength{\intextsep}{6pt plus 1.0pt minus 1.5pt}

\begin{document}

\title{Parallel Expectation Maximization Clustering\\
{\normalsize\normalfont\itshape Department of Information Engineering and Computer Science, University of Trento}\\
{\normalsize\normalfont\url{https://github.com/HPC-2627/ParallelExpectationMaximizationClustering}}}

\author{
\IEEEauthorblockN{Michael Bernasconi}
\IEEEauthorblockA{\textit{Student ID: 267681}\\
\textit{michael.bernasconi@studenti.unitn.it}}
\and
\IEEEauthorblockN{Gabriele Fronzoni}
\IEEEauthorblockA{\textit{Student ID: 267336}\\
\textit{gabriele.fronzoni@studenti.unitn.it}}
}

\maketitle

\begin{abstract}
\end{abstract}

\begin{IEEEkeywords}
\end{IEEEkeywords}

\section{Introduction}
Il clustering raggruppa i punti di un dataset in modo che quelli dello stesso gruppo siano simili. Qui si assume che i dati provengano da una mistura di distribuzioni gaussiane e che il gruppo di ogni punto non sia noto; i parametri della mistura si stimano con l'algoritmo Expectation--Maximization (EM). EM è iterativo e ogni iterazione scorre tutto il dataset: il costo cresce con i dati ed è un buon candidato alla parallelizzazione. L'obiettivo del progetto è implementare EM in C con un design ibrido MPI+OpenMP, valutare ottimizzazioni di calcolo e comunicazione e misurarne la scalabilità su più nodi (strong e weak scaling) su dataset reali e sintetici.

\subsection{L'algoritmo EM per misture di gaussiane}
Distinguiamo due cose: il \emph{modello} è l'ipotesi sulla forma dei dati (una mistura di gaussiane, con parametri ancora ignoti), mentre \emph{EM} è il procedimento iterativo che stima quei parametri dai dati.

\paragraph{Modello.} L'input è una matrice $X$ di $N$ punti (righe) e $D$ variabili (colonne); $x_i\in\mathbb{R}^D$ è la riga $i$. Il numero di cluster $K$ è scelto dall'utente e non si usano etichette. Si assume che i punti provengano da una mistura di $K$ gaussiane:
\begin{equation}
p(x)=\sum_{k=1}^{K}\pi_k\,\mathcal{N}(x\mid\mu_k,\sigma_k^2).
\end{equation}
Per il cluster $k$: $\pi_k$ è il \emph{peso} (frazione di punti, $\sum_k\pi_k=1$), $\mu_k\in\mathbb{R}^D$ la \emph{media} (centro) e $\sigma^2_k=(\sigma^2_{k1},\dots,\sigma^2_{kD})$ le \emph{varianze} delle $D$ variabili. Usiamo la covarianza \emph{diagonale}: dentro un cluster le variabili sono indipendenti, quindi
\begin{equation}
\ln\mathcal{N}(x\mid\mu_k,\sigma_k^2)=-\tfrac12\Big[D\ln 2\pi+\sum_{d=1}^{D}\ln\sigma^2_{kd}+m_k(x)\Big],
\end{equation}
\begin{equation}
m_k(x)=\sum_{d=1}^{D}\frac{(x_d-\mu_{kd})^2}{\sigma^2_{kd}},
\end{equation}
dove $x_d$ e $\mu_{kd}$ sono le componenti $d$ di $x$ e $\mu_k$, e $m_k(x)$ misura quanto $x$ è lontano dal centro del cluster (in unità di deviazione standard). Si cercano i parametri $\theta=\{\pi_k,\mu_k,\sigma^2_k\}_{k=1}^{K}$ che massimizzano la log-verosimiglianza, cioè quanto bene il modello spiega i dati:
\begin{equation}
\mathcal{L}(\theta)=\sum_{i=1}^{N}\ln\sum_{k=1}^{K}\pi_k\,\mathcal{N}(x_i\mid\mu_k,\sigma_k^2).
\end{equation}
Poiché il cluster di ogni punto non è noto, non c'è una formula diretta: EM~\cite[Sez.~9.2.2]{bishop} parte da una stima iniziale e alterna due passi (Algoritmo~\ref{alg:em}).

\paragraph{E-step.} Per ogni punto $i$ e cluster $k$ calcola la \emph{responsabilità} $r_{ik}\in[0,1]$, la probabilità che $x_i$ appartenga al cluster $k$ con i parametri correnti ($\sum_k r_{ik}=1$); $j$ scorre tutti i cluster:
\begin{equation}
r_{ik}=\frac{\pi_k\,\mathcal{N}(x_i\mid\mu_k,\sigma_k^2)}{\sum_{j=1}^{K}\pi_j\,\mathcal{N}(x_i\mid\mu_j,\sigma_j^2)}.
\end{equation}

\paragraph{M-step.} Ricalcola i parametri pesando ogni punto con $r_{ik}$. Con $N_k=\sum_{i=1}^{N}r_{ik}$ (numero effettivo di punti del cluster $k$):
\begin{gather}
\mu_k=\tfrac{1}{N_k}\sum_{i=1}^{N} r_{ik}\,x_i,\qquad
\pi_k=\tfrac{N_k}{N},\\
\sigma^2_{kd}=\tfrac{1}{N_k}\sum_{i=1}^{N} r_{ik}\,x_{id}^2-\mu_{kd}^2.
\end{gather}
Nel codice l'M-step accumula in un solo passaggio le somme $N_k$, $\sum_i r_{ik}x_i$ e $\sum_i r_{ik}x_i^2$ (quest'ultima componente per componente) e poi le converte nei parametri.

\paragraph{Inizializzazione e arresto.} Le medie $\mu_k$ sono $K$ righe distinte di $X$ scelte con un generatore con seed (così ogni processo ottiene le stesse righe senza comunicare), $\pi_k=1/K$ e $\sigma^2_{kd}=1$ (i dati sono standardizzati: media 0 e varianza 1 per variabile). Dopo ogni E-step si calcola $\mathcal{L}$ e ci si ferma se $|\mathcal{L}-\mathcal{L}_{\mathrm{prev}}|\le\varepsilon|\mathcal{L}|$, con $\mathcal{L}_{\mathrm{prev}}$ il valore dell'iterazione precedente ($-\infty$ all'inizio) ed $\varepsilon$ una tolleranza relativa, oppure dopo $T_{\max}$ iterazioni (default $\varepsilon=10^{-6}$, $T_{\max}=200$). $\mathcal{L}$ non decresce mai~\cite[Sez.~9.4]{bishop} ma si raggiunge solo un massimo locale: con $R$ restart si tiene l'esecuzione con $\mathcal{L}$ più alta (nelle misure $R=1$).

\begin{algorithm}[t]
\caption{EM per una mistura di gaussiane (una esecuzione)}\label{alg:em}
\small
\KwIn{$X$ ($N\times D$), $K$, $\varepsilon$, $T_{\max}$}
\KwOut{$\theta$, etichette $\arg\max_k r_{ik}$}
Inizializza $\theta$: $\mu_k\leftarrow$ righe distinte casuali, $\pi_k\leftarrow 1/K$, $\sigma^2_{kd}\leftarrow 1$\;
$\mathcal{L}_{\mathrm{prev}}\leftarrow-\infty$\;
\For{$t=1$ \KwTo $T_{\max}$}{
  \textbf{E-step:} calcola $r_{ik}$ per ogni $i,k$ e $\mathcal{L}$\;
  \lIf{$|\mathcal{L}-\mathcal{L}_{\mathrm{prev}}|\le\varepsilon|\mathcal{L}|$}{\textbf{break}}
  \textbf{M-step:} calcola $N_k$, $\mu_k$, $\pi_k$, $\sigma^2_{kd}$ per ogni $k,d$\;
  $\mathcal{L}_{\mathrm{prev}}\leftarrow\mathcal{L}$\;
}
\lIf{$T_{\max}$ raggiunto senza convergenza}{ripeti l'E-step (per $\mathcal{L}$ e $r_{ik}$ coerenti con $\theta$)}
\end{algorithm}

\paragraph{Accorgimenti numerici.} Le formule, implementate alla lettera, producono errori in aritmetica a precisione finita; il codice adotta tre protezioni.
\begin{itemize}
\item \textbf{Calcolo in scala logaritmica.} Con $D$ grande le densità $\mathcal{N}$ sono numeri minuscoli che il calcolatore arrotonda a 0, e la normalizzazione di $r_{ik}$ darebbe $0/0$. Per ogni punto si calcolano quindi i logaritmi $v_k=\ln\big(\pi_k\,\mathcal{N}(x_i\mid\mu_k,\sigma_k^2)\big)$ e li si somma con
$\ln\sum_{k}e^{v_k}=m+\ln\sum_{k}e^{v_k-m}$, dove $m=\max_k v_k$: il termine più grande diventa $e^0=1$ e la somma non è mai nulla. Il risultato è $\ln p(x_i)$, che contribuisce a $\mathcal{L}$ e dà $r_{ik}=e^{v_k-\ln p(x_i)}$.
\item \textbf{Varianze non nulle.} Se un cluster si riduce a pochi punti quasi identici, $\sigma^2_{kd}$ tende a 0 e la densità a infinito, rendendo $\mathcal{L}$ instabile~\cite[Sez.~9.2.1]{bishop}. Si impone perciò $\sigma^2_{kd}\ge\delta=10^{-6}$ a ogni M-step.
\item \textbf{Cluster vuoti.} Se $N_k<10^{-10}$ nessun punto è assegnato al cluster e le formule dividerebbero per zero: media e varianze restano quelle precedenti e il peso viene portato al minimo $\pi_k=10^{-12}$, mai 0, perché l'E-step usa $\ln\pi_k$.
\end{itemize}

\subsection{Implementazioni sequenziali e costo}
Abbiamo scritto due versioni \emph{sequenziali}, \emph{solo per capire a fondo l'algoritmo}: non sono la baseline delle prestazioni, che è invece la versione MPI+OpenMP diagonale, rispetto alla quale si valutano tutte le ottimizzazioni (descritte nelle sezioni successive).
\begin{itemize}
\item \textbf{Diagonale:} è il modello descritto sopra, con una varianza $\sigma^2_{kd}$ per variabile.
\item \textbf{Piena:} ogni cluster ha una matrice di covarianza completa $\Sigma_k\in\mathbb{R}^{D\times D}$, che modella anche le correlazioni tra variabili. Il modello cambia, quindi i suoi risultati si confrontano solo con altre esecuzioni a covarianza piena.
\end{itemize}

\paragraph{Costo di un'iterazione.} Con $N$ punti, $K$ cluster e $D$ variabili:
\begin{itemize}
\item \textbf{Diagonale, $O(NKD)$.} Per ognuna delle $NK$ coppie $(i,k)$ l'E-step calcola $m_k(x_i)$ sommando $D$ termini, e l'M-step accumula $D$ termini in $\sum_i r_{ik}x_i$ e $D$ in $\sum_i r_{ik}x_i^2$.
\item \textbf{Piena, $O(NKD^2)$.} Per ogni coppia $(i,k)$ la distanza di Mahalanobis e l'accumulo di $\Sigma_k$ costano $O(D^2)$ invece di $O(D)$. Si aggiungono $O(KD^3)$ per fattorizzare le $K$ matrici $\Sigma_k$; non dipende da $N$ ed è trascurabile se $N\gg D$.
\item \textbf{Comune, $O(NK)$.} La normalizzazione delle responsabilità e il calcolo di $\mathcal{L}$ costano $K$ operazioni per punto.
\end{itemize}


\section{Parallel Design Methodology}\label{sec:design}
\subsection{Foster's Design Strategy}
\subsection{Communication}

\section{Experimental Setup \& Dataset}
\subsection{Hardware Infrastructure}
\subsection{Dataset}

\begin{table*}[t]
\centering
\caption{Real datasets. $N$ and $D$ are the values actually used by the EM after preprocessing; $K$ is the number of ground-truth classes (used only for evaluation).}
\label{tab:real}
\footnotesize
\setlength{\tabcolsep}{3pt}
\begin{tabular}{llrrrrll}
\toprule
\textbf{Name} & \textbf{Source} & \textbf{$N$} & \textbf{$D$} & \textbf{$K$} & \textbf{File} & \textbf{Label} & \textbf{Description} \\
\midrule
\texttt{statlog\_heart}   & \srclink{https://archive.ics.uci.edu/dataset/145}{UCI \#145} & 270 & 13 & 2 & 28 KB & \texttt{heart-disease} & Heart disease, small clinical sample \\
\texttt{heart\_failure}   & \srclink{https://archive.ics.uci.edu/dataset/519}{UCI \#519} & 299 & 12 & 2 & 29 KB & \texttt{death\_event} & Heart failure patients, survival \\
\texttt{gene\_rnaseq}     & \srclink{https://archive.ics.uci.edu/dataset/401}{UCI \#401} & 801 & 20\,264 & 5 & 130 MB & \texttt{Class} & TCGA RNA-Seq, 5 cancer types \\
\texttt{cardiotocography} & \srclink{https://archive.ics.uci.edu/dataset/193}{UCI \#193} & 2\,126 & 21 & 10 & 357 KB & \texttt{CLASS} & Fetal cardiotocograms \\
\texttt{parkinson}        & \srclink{https://archive.ics.uci.edu/dataset/189}{UCI \#189} & 5\,875 & 19 & -- & 893 KB & -- (regression) & Voice telemonitoring \\
\texttt{eeg\_eye\_state}  & \srclink{https://archive.ics.uci.edu/dataset/264}{UCI \#264} & 14\,980 & 14 & 2 & 1.7 MB & \texttt{eyeDetection} & EEG, eyes open/closed \\
\texttt{cardiovascular}   & \srclink{https://www.openml.org/d/45547}{OpenML 45547} & 70\,000 & 11 & 2 & 6.2 MB & \texttt{cardio} & Cardiovascular disease \\
\texttt{diabetes130}      & \srclink{https://archive.ics.uci.edu/dataset/296}{UCI \#296} & 101\,766 & 11 & 3 & 9.0 MB & \texttt{readmitted} & Hospital readmission, 130 US hospitals \\
\texttt{cdc\_diabetes}    & \srclink{https://archive.ics.uci.edu/dataset/891}{UCI \#891} & 253\,680 & 21 & 2 & 43 MB & \texttt{Diabetes\_binary} & CDC health indicators survey \\
\texttt{brainweb}         & \srclink{https://brainweb.bic.mni.mcgill.ca/}{BrainWeb (subj.~44)} & 11\,862\,016 & 1 & 12 & 95 MB & crisp labels & Simulated T1 brain MRI, one voxel = one point \\
\bottomrule
\end{tabular}
\end{table*}

\begin{table*}[t]
\centering
\caption{Synthetic profiles. $\delta$ is the distance between cluster means in units of the within-cluster standard deviation.}
\label{tab:synth-profiles}
\footnotesize
\setlength{\tabcolsep}{3pt}
\begin{tabular}{llll}
\toprule
 & \texttt{synth\_clean} & \texttt{synth\_overlap} & \texttt{synth\_rare} \\
\midrule
Role & easy, satisfies the model & realistic & hard \\
Clusters $K$ & 4 & 5 & 6 \\
Variables $D$ & 10 & 20 & 50 \\
Cluster weights & 25, 25, 25, 25\,\% & 30, 25, 20, 15, 10\,\% & 40, 25, 18, 12, 4, 1\,\% \\
Separation $\delta$ & 6.0 (well separated) & 2.5 (overlapping) & 4.0 (intermediate) \\
Covariance in a cluster & diagonal & correlated (3-factor model, 0.5) & diagonal \\
Std.\ per variable & 1.0 & 1.0 (before correlation) & in $[0.6, 1.4]$ \\
Outliers & 0\,\% & 0\,\% & 2\,\% (label $-1$) \\
Skewed (log-normal) variables & no & yes (7 of 20) & no \\
Physiological clipping & no & yes & yes (clinical variables) \\
What it tests & correctness, scaling reference & slow convergence, choice of $K$, model mismatch & numerical stability \\
\bottomrule
\end{tabular}
\end{table*}

\begin{table*}[t]
\centering
\caption{Clinical variables used to build the synthetic data (typical mean $\pm$ std and physiological limits).}
\label{tab:synth-vars}
\footnotesize
\begin{tabular}{rlllcc}
\toprule
\# & \textbf{Variable} & \textbf{Unit} & \textbf{Mean $\pm$ std} & \textbf{Limits} & \textbf{Skew} \\
\midrule
1 & \texttt{age} & years & $55 \pm 15$ & 18--100 & \\
2 & \texttt{bmi} & kg/m$^2$ & $27 \pm 5$ & 14--60 & \\
3 & \texttt{sys\_bp} & mmHg & $130 \pm 20$ & 80--230 & \\
4 & \texttt{dia\_bp} & mmHg & $80 \pm 11$ & 40--130 & \\
5 & \texttt{glucose} & mg/dL & $105 \pm 30$ & 50--400 & \checkmark \\
6 & \texttt{hba1c} & \% & $5.8 \pm 1.0$ & 4--14 & \checkmark \\
7 & \texttt{cholesterol} & mg/dL & $200 \pm 40$ & 90--400 & \\
8 & \texttt{hdl} & mg/dL & $52 \pm 14$ & 20--110 & \\
9 & \texttt{creatinine} & mg/dL & $1.0 \pm 0.3$ & 0.3--6 & \checkmark \\
10 & \texttt{heart\_rate} & bpm & $75 \pm 12$ & 40--150 & \\
11 & \texttt{ldl} & mg/dL & $120 \pm 35$ & 30--300 & \\
12 & \texttt{triglycerides} & mg/dL & $140 \pm 70$ & 30--800 & \checkmark \\
13 & \texttt{albumin} & g/dL & $4.2 \pm 0.4$ & 2--5.5 & \\
14 & \texttt{hemoglobin} & g/dL & $13.5 \pm 1.8$ & 6--19 & \\
15 & \texttt{wbc} & $10^3/\mu$L & $7.0 \pm 2.0$ & 2--25 & \checkmark \\
16 & \texttt{platelets} & $10^3/\mu$L & $250 \pm 60$ & 50--600 & \\
17 & \texttt{sodium} & mmol/L & $140 \pm 3$ & 120--160 & \\
18 & \texttt{potassium} & mmol/L & $4.2 \pm 0.4$ & 2.5--6.5 & \\
19 & \texttt{alt} & U/L & $28 \pm 15$ & 5--200 & \checkmark \\
20 & \texttt{crp} & mg/L & $3.0 \pm 3.0$ & 0.1--100 & \checkmark \\
21--50 & \texttt{biomarker\_01..30} & -- & $0 \pm 1$ & none & \\
\bottomrule
\end{tabular}
\end{table*}

\begin{table*}[t]
\centering
\caption{Weak scaling: rows and data file size for each profile ($N_0 = 250\,000$ rows per process; labels file of $4N$ bytes for every profile). MB $= 10^6$ bytes.}
\label{tab:weak}
\footnotesize
\setlength{\tabcolsep}{3pt}
\begin{tabular}{rrrrrr}
\toprule
\textbf{$p$} & \textbf{$N = N_0 \cdot p$} & \texttt{synth\_clean} & \texttt{synth\_overlap} & \texttt{synth\_rare} & \textbf{Labels} \\
\midrule
1  & 250\,000   & 20 MB  & 40 MB  & 100 MB   & 1 MB \\
2  & 500\,000   & 40 MB  & 80 MB  & 200 MB   & 2 MB \\
4  & 1\,000\,000 & 80 MB  & 160 MB & 400 MB   & 4 MB \\
8  & 2\,000\,000 & 160 MB & 320 MB & 800 MB   & 8 MB \\
16 & 4\,000\,000 & 320 MB & 640 MB & 1\,600 MB & 16 MB \\
\midrule
\textbf{Total} & & 620 MB & 1\,240 MB & 3\,100 MB & 31 MB \\
\bottomrule
\end{tabular}
\end{table*}

\begin{table*}[t]
\centering
\caption{Strong scaling (fixed $N = 4\,000\,000$): data per rank and work per rank $N \cdot K \cdot D / p$. The work is an operation count proportional to the cost of an iteration, not a measured time.}
\label{tab:strong}
\footnotesize
\setlength{\tabcolsep}{3pt}
\begin{tabular}{rrrrrrrr}
\toprule
 & & \multicolumn{3}{c}{\textbf{Data per rank}} & \multicolumn{3}{c}{\textbf{Work per rank ($N K D / p$)}} \\
\cmidrule(lr){3-5} \cmidrule(lr){6-8}
\textbf{$p$} & \textbf{Rows/rank} & \texttt{clean} & \texttt{overlap} & \texttt{rare} & \texttt{clean} & \texttt{overlap} & \texttt{rare} \\
\midrule
1  & 4\,000\,000 & 320 MB & 640 MB & 1\,600 MB & $1.6 \cdot 10^8$ & $4.0 \cdot 10^8$ & $1.2 \cdot 10^9$ \\
2  & 2\,000\,000 & 160 MB & 320 MB & 800 MB    & $8.0 \cdot 10^7$ & $2.0 \cdot 10^8$ & $6.0 \cdot 10^8$ \\
4  & 1\,000\,000 & 80 MB  & 160 MB & 400 MB    & $4.0 \cdot 10^7$ & $1.0 \cdot 10^8$ & $3.0 \cdot 10^8$ \\
8  & 500\,000    & 40 MB  & 80 MB  & 200 MB    & $2.0 \cdot 10^7$ & $5.0 \cdot 10^7$ & $1.5 \cdot 10^8$ \\
16 & 250\,000    & 20 MB  & 40 MB  & 100 MB    & $1.0 \cdot 10^7$ & $2.5 \cdot 10^7$ & $7.5 \cdot 10^7$ \\
\bottomrule
\end{tabular}
\end{table*}

\section{Performance Evaluation}
\subsection{Execution Time and Throughput}
\subsection{Strong Scaling}
\subsection{Weak Scaling}
\subsection{Communication vs Computation Breakdown}

\section{Discussion}

\section{Conclusion}

\clearpage
\begin{thebibliography}{1}
\bibitem{bishop} C. M. Bishop, \href{https://www.microsoft.com/en-us/research/uploads/prod/2006/01/Bishop-Pattern-Recognition-and-Machine-Learning-2006.pdf}{\textcolor{blue}{\emph{Pattern Recognition and Machine Learning}}}, Springer, 2006, ch.~9.
\end{thebibliography}

\end{document}
